\section{Supplementary Corollaries S2.6--S2.7: Selection-Safe Certification}

This appendix formalizes two selection-validity extensions of the finite-sample certificate.

\subsection{S2.6: arbitrary independent training}

Let $T$ denote an arbitrary training random element. Conditional on $T$, suppose a training procedure outputs a predictor/accessibility rule
\[
Y_T=f_T(X),
\qquad
S_T=s_T(Y_T),
\]
and bounds
\[
|Y_T|\le B_Y(T),
\qquad
0\le S_T\le B_S(T),
\qquad
|U-Y_T|\le B_R(T)
\]
for fresh population draws, almost surely. Let the certification observations be i.i.d. from the target population and independent of $T$.

For the trained rule, define
\[
C(T)=\operatorname{Cov}(U,S_T\mid T),
\]
and construct the S2.5 lower certificate $D_L(T)$ at confidence level $\delta$ from the independent certification sample.

\begin{corollary}[Selection validity after independent training]
Under the assumptions above,
\[
P\!\left(C(T)\ge D_L(T)\mid T\right)\ge1-\delta
\]
almost surely. Consequently,
\[
P\!\left(C(T)\ge D_L(T)\right)\ge1-\delta.
\]
\end{corollary}

\begin{proof}
Fix a realized training state $T=t$. Conditional on $T=t$, the predictor, accessibility rule, and bounds are fixed before the independent certification observations are evaluated. Therefore S2.5 applies conditionally and gives
\[
P\!\left(C(t)\ge D_L(t)\mid T=t\right)\ge1-\delta.
\]
Taking expectation over $T$ yields
\[
P(C(T)\ge D_L(T))
=
\mathbb E\!\left[P(C(T)\ge D_L(T)\mid T)\right]
\ge1-\delta.
\]
\end{proof}

Thus the complexity of the upstream learning procedure is not itself a validity problem when the final certification sample remains independent.

\subsection{S2.7: finite candidate selection with family-wise correction}

Suppose a training stage independent of the certification sample produces $K$ candidate rules
\[
(Y^{(k)},S^{(k)}),
\qquad
k=1,\ldots,K,
\]
with candidate definitions and valid bounds fixed before the certification observations are inspected.

For each candidate, let
\[
C_k=\operatorname{Cov}(U,S^{(k)}).
\]
Apply S2.5 to candidate $k$ at confidence level
\[
\delta_k=\frac{\delta}{K}.
\]
With equal allocation, the corresponding Hoeffding radius is
\[
\tau_{n,\delta,K}
=
\sqrt{\frac{\log(10K/\delta)}{2n}}.
\]
Let $D_{L,k}$ denote the resulting lower certificate.

\begin{corollary}[Finite-candidate post-selection validity]
The candidate-specific certificates satisfy
\[
P\!\left(
C_k\ge D_{L,k}
\text{ for all }k=1,\ldots,K
\right)
\ge1-\delta.
\]
Hence, for any data-dependent selection rule
\[
\widehat k
\in\{1,\ldots,K\}
\]
computed from the certification sample,
\[
P\!\left(
C_{\widehat k}
\ge
D_{L,\widehat k}
\right)
\ge1-\delta.
\]
\end{corollary}

\begin{proof}
For each fixed $k$, S2.5 gives
\[
P(C_k<D_{L,k})\le\frac{\delta}{K}.
\]
Therefore the union bound gives
\[
P\!\left(\exists k: C_k<D_{L,k}\right)
\le
\sum_{k=1}^K\frac{\delta}{K}
=
\delta.
\]
On the complementary simultaneous-validity event, every candidate lower bound is valid. Any index selected from the certification sample therefore satisfies
\[
C_{\widehat k}\ge D_{L,\widehat k}.
\]
\end{proof}

The same argument permits unequal allocations $\delta_k=\delta w_k$ for positive weights satisfying $\sum_k w_k\le1$.

\subsection{Boundary of the selection results}

The finite-family result does not justify uncorrected best-of-$K$ reporting. Nor does it cover predictor forms, accessibility rules, clipping bounds, or new candidate definitions invented after observing the certification sample. Those settings require a new independent sample or a more general selective-inference or uniform-convergence argument.

S2.6--S2.7 are statistical-validity results only. They preserve the abstract covariance certificate after realistic training and finite candidate selection; they do not supply a physical Everett accessibility law.
